\documentclass[aspectratio=169,11pt]{beamer}

\usepackage{fontspec}
\usepackage[russian]{babel}
\setmainfont{DejaVu Sans}
\setsansfont{DejaVu Sans}
\setmonofont{DejaVu Sans Mono}

\usepackage{amsmath,mathtools}
\usepackage{booktabs,tabularx,array,multirow}
\usepackage{tikz}
\usepackage{pgfplots}
\usepackage{ragged2e}
\usetikzlibrary{arrows.meta,positioning,calc,shapes.misc}
\pgfplotsset{compat=1.18}

% --- style ---
\definecolor{msuBlue}{HTML}{123A63}
\definecolor{textDark}{HTML}{1B1B1B}
\definecolor{softBlue}{HTML}{EAF1F8}
\definecolor{lineGray}{HTML}{D9E1E8}
\definecolor{mutedText}{HTML}{52616F}
\definecolor{pointBlue}{HTML}{1F77B4}
\definecolor{goodGreen}{HTML}{2E7D32}

\setbeamertemplate{navigation symbols}{}
\setbeamertemplate{blocks}[rounded][shadow=false]
\setbeamertemplate{itemize item}{\color{msuBlue}$\blacktriangleright$}
\setbeamertemplate{itemize subitem}{\color{msuBlue}$\circ$}
\setbeamertemplate{footline}{}
\setbeamertemplate{headline}{\vspace{1mm}\hspace{5mm}\color{lineGray}\rule{\dimexpr\paperwidth-10mm}
{0.4pt}\vspace{-1mm}}

\setbeamercolor{normal text}{fg=textDark,bg=white}
\setbeamercolor{frametitle}{fg=msuBlue,bg=white}
\setbeamercolor{title}{fg=msuBlue}
\setbeamercolor{block title}{fg=white,bg=msuBlue}
\setbeamercolor{block body}{fg=textDark,bg=softBlue}
\setbeamerfont{frametitle}{size=\Large,series=\bfseries}
\setbeamerfont{title}{size=\Huge,series=\bfseries}
\setbeamersize{text margin left=8mm,text margin right=8mm}

\newcommand{\tight}{\vspace{-1mm}}
\newcommand{\term}[1]{\textbf{\color{msuBlue}#1}}
\newcommand{\smallnote}[1]{{\color{mutedText}\small #1}}
\newcolumntype{Y}{>{\RaggedRight\arraybackslash}X}
\newcolumntype{C}[1]{>{\centering\arraybackslash}p{#1}}

\title{Поиск оптимального критерия остановки итеративной процедуры}
\author{Якунин Владимир Петрович}
\date{2026}

\begin{document}

% Числовые блоки синхронизированы с tex/tables/*.csv через scripts/render_presentation_data.py
% AUTO-GENERATED — do not edit by hand.
% Source: tex/tables/wikitext_lm_base.csv, arc_stopping_base.csv,
%          wikitext_lm_ft_compare.csv, arc_stopping_finetune_compare.csv
% Regenerate: uv run python scripts/render_presentation_data.py

\providecommand{\PresWikiTabBody}{%
Конечное дискретное & 0,90 & 4,73 & 0,743 & 0,784 \\
Сглаженная разметка & 0,60 & 5,00 & 0,744 & 2,083 \\
Смесь геометрических & 0,90 & 3,12 & 0,743 & 0,825 \\
Смесь экспоненциальных & 0,90 & 3,06 & 0,743 & 0,822 \\
Степенной спад & 0,90 & 4,66 & 0,743 & 0,843 \\
}

\providecommand{\PresArcTabBody}{%
Конечное дискретное & 0,90 & 3,44 & 0,730 & 1,126 \\
Сглаженная разметка & 0,60 & 5,00 & 0,737 & 2,015 \\
Смесь геометрических & 0,90 & 4,25 & 0,745 & 0,974 \\
Смесь экспоненциальных & 0,80 & 3,50 & 0,742 & 1,099 \\
Степенной спад & 0,80 & 3,71 & 0,741 & 1,125 \\
Отрицательное биномиальное & 0,90 & 3,85 & 0,731 & 1,271 \\
}

\providecommand{\PresWikiQcXmin}{2.9061}
\providecommand{\PresWikiQcXmax}{5.1551}
\providecommand{\PresWikiQcYmin}{0.742898}
\providecommand{\PresWikiQcYmax}{0.744071}

\providecommand{\PresArcQcXmin}{3.3152}
\providecommand{\PresArcQcXmax}{5.1248}
\providecommand{\PresArcQcYmin}{0.724058}
\providecommand{\PresArcQcYmax}{0.750550}

\providecommand{\PresWikiScatter}{%
\addplot coordinates {(4.7331,0.743315)};
\addlegendentry{Конечное дискретное}
\addplot coordinates {(5,0.743721)};
\addlegendentry{Сглаженная разметка}
\addplot coordinates {(3.1172,0.743362)};
\addlegendentry{Смесь геометрических}
\addplot coordinates {(3.0612,0.743312)};
\addlegendentry{Смесь экспоненциальных}
\addplot coordinates {(4.6641,0.743248)};
\addlegendentry{Степенной спад}
}

\providecommand{\PresArcScatter}{%
\addplot coordinates {(3.44,0.729512)};
\addlegendentry{Конечное дискретное}
\addplot coordinates {(5,0.737145)};
\addlegendentry{Сглаженная разметка}
\addplot coordinates {(4.2475,0.745096)};
\addlegendentry{Смесь геометрических}
\addplot coordinates {(3.4975,0.741569)};
\addlegendentry{Смесь экспоненциальных}
\addplot coordinates {(3.71,0.741110)};
\addlegendentry{Степенной спад}
\addplot coordinates {(3.85,0.731179)};
\addlegendentry{Отрицательное биномиальное}
}

\providecommand{\PresFTtauBody}{%
Конечное дискретное & 4,91 & 4,73 & 3,44 & 1,98 \\
Сглаженная разметка & 8,00 & 5,00 & 5,00 & 5,00 \\
Смесь геометрических & 8,00 & 3,12 & 4,25 & 1,74 \\
Смесь экспоненциальных & 8,00 & 3,06 & 3,50 & 3,44 \\
Степенной спад & 8,00 & 4,66 & 3,71 & 1,87 \\
Отрицательное биномиальное & --- & --- & 3,85 & 2,35 \\
}

\providecommand{\PresWikiTabConclusion}{Максимальный топ-1 на контроле после дополнительного обучения головы оракула: \textbf{Сглаженная разметка} (0,744). Различия между семьями малы; подробнее --- раздел WikiText в диссертации.}

\providecommand{\PresArcTabConclusion}{Максимальная токеновая точность (ARC): \textbf{Смесь геометрических} (0,745). Смесь экспоненциальных даёт сопоставимое качество при меньшем $\tau$.}



% 1
\begin{frame}[plain]
\centering
\vspace{1mm}
{\color{msuBlue}\rule{0.80\paperwidth}{0.7pt}}\par
\vspace{5mm}
{\color{msuBlue}\bfseries\fontsize{21}{25}\selectfont Поиск оптимального критерия\\[1mm] остановки
итеративной процедуры}\par
\vspace{4mm}
{\color{msuBlue}\rule{16mm}{1.0pt}}\par
\vspace{4mm}
{\large Защита выпускной квалификационной работы}\par
\vspace{6mm}
{\normalsize\bfseries Якунин Владимир Петрович}\par
\vspace{1.5mm}
{\normalsize Научный руководитель: Хохлов Юрий Степанович}\par
\vspace{5mm}
{\small МГУ имени М. В. Ломоносова, ВМК, кафедра математической статистики}\par
\vfill
{\normalsize 2026}\par
\vspace{1mm}
\end{frame}

% 2
\begin{frame}{1. Задача в простых словах}
\begin{block}{Что нужно сделать}
Для каждой позиции текста выбрать, сколько внутренних уточнений выполнить перед выдачей прогноза
следующего токена.
\end{block}

\vspace{2mm}
\begin{columns}[T,onlytextwidth]
\begin{column}{0.53\textwidth}
\small
\term{Оракул} — лучший шаг, выбранный постфактум: известны правильный ответ и потери на всех
шагах.\par\vspace{2mm}
\term{Время остановки} — номер внутреннего шага, на котором рекурсивная процедура
прекращается.\par\vspace{2mm}
\term{Оптимальность} — минимум суммы ошибки прогноза и штрафа за вычисления.
\end{column}
\begin{column}{0.42\textwidth}
\begin{tikzpicture}[node distance=5mm, every node/.style={font=\small}]
\node[draw=msuBlue,rounded corners,fill=softBlue,minimum width=40mm,minimum height=7mm] (a)
{префикс текста};
\node[draw=msuBlue,rounded corners,fill=softBlue,minimum width=40mm,minimum height=7mm,below=of a]
(b) {внутренние шаги};
\node[draw=msuBlue,rounded corners,fill=softBlue,minimum width=40mm,minimum height=7mm,below=of b]
(c) {распределение времени};
\node[draw=msuBlue,rounded corners,fill=softBlue,minimum width=40mm,minimum height=7mm,below=of c]
(d) {правило остановки};
\draw[-{Latex[length=2mm]},thick,color=msuBlue] (a)--(b);
\draw[-{Latex[length=2mm]},thick,color=msuBlue] (b)--(c);
\draw[-{Latex[length=2mm]},thick,color=msuBlue] (c)--(d);
\end{tikzpicture}
\end{column}
\end{columns}
\end{frame}

% 3
\begin{frame}{2. Что было до этой работы}
\small
\small\begin{tabularx}{\textwidth}{@{}p{0.31\textwidth}Y@{}}
\toprule
\textbf{Направление} & \textbf{Что дает для постановки} \\
\midrule
Последовательный анализ и оптимальная остановка & язык риска, остановки и сравнения правил: Вальд,
Снелл, Ширяев.\\[1mm]
Статистическая теория решений & разделяет прогноз неизвестной величины и действие, принимаемое по
прогнозу.\\[1mm]
Адаптивная глубина и ранний выход нейросетей & показывают, что вычисления можно тратить
неравномерно, но часто используют эвристику уверенности.\\[1mm]
Малые рекурсивные модели & дают механизм многократного уточнения скрытого состояния одним
рекурсивным ядром.\\
\bottomrule
\end{tabularx}

\vspace{4mm}
\vspace{2mm}
\smallnote{Разрыв: вместо прямой эвристики остановки оценивается полный закон оракульного времени,
а правило строится поверх него.}
\end{frame}

% 4
\begin{frame}{3. Основная идея}
\begin{columns}[T,onlytextwidth]
\begin{column}{0.53\textwidth}
\begin{block}{Оракульная метка}
\vspace{-1mm}
\[
L_{t,k}=-\log Q_{t,k}(X_{t+1}),
\]
\vspace{-3mm}
\[
\tau_t^\star=\arg\min_{0\le j\le K}\{L_{t,j}+\lambda c(j)\}.
\]
\vspace{-1mm}
\smallnote{Метка вычисляется только при обучении: нужен правильный следующий токен и вся
траектория шагов.}
\end{block}
\end{column}
\begin{column}{0.43\textwidth}
\begin{block}{Модель времени}
\vspace{-1mm}
\[
\widehat p_{t,k}(j)=P_\theta(\tau_t^\star=j\mid H_{t,k}).
\]
\vspace{-1mm}
На каждом текущем шаге $k$ модель выдает вероятности всех возможных оракульных шагов $j=0,\ldots,K$.
\end{block}
\end{column}
\end{columns}

\vspace{2mm}
\centering
\begin{tikzpicture}[every node/.style={font=\small}, node distance=4mm]
\node[draw=msuBlue,fill=softBlue,rounded corners,minimum width=32mm,minimum height=8mm] (p)
{прошлое: $j<k$};
\node[draw=msuBlue,fill=softBlue,rounded corners,minimum width=32mm,minimum height=8mm,right=of p]
(n) {сейчас: $j=k$};
\node[draw=msuBlue,fill=softBlue,rounded corners,minimum width=32mm,minimum height=8mm,right=of n]
(f) {будущее: $j>k$};
\end{tikzpicture}
\end{frame}

% 5
\begin{frame}{4. Что конкретно сделано}
\begin{columns}[T,onlytextwidth]
\begin{column}{0.48\textwidth}
\begin{block}{Модель}
\begin{itemize}\setlength\itemsep{1.5mm}
\item рекурсивное уточнение распределения $Q_{t,k}$;
\item голова времени остановки для $\widehat p_{t,k}$;
\item разделение прогноза распределения и правила действия.
\end{itemize}
\end{block}
\end{column}
\begin{column}{0.48\textwidth}
\begin{block}{Эксперименты}
\begin{itemize}\setlength\itemsep{1.5mm}
\item шесть семейств распределений на $\{0,\ldots,K\}$;
\item пять правил остановки;
\item проверка на WikiText-103 и ARC-AGI;
\item сравнение базового обучения и дообучения головы оракула.
\end{itemize}
\end{block}
\end{column}
\end{columns}

\vspace{3mm}
\smallnote{Преимущество: правило использует массу вероятности на прошлом, текущем и будущем
оптимуме, а не одну эвристику уверенности.}
\end{frame}

% 6
\begin{frame}{5. Семейства распределений и правила остановки}
\small
\begin{columns}[T,onlytextwidth]
\begin{column}{0.50\textwidth}
\textbf{Семейства распределений}
\vspace{1mm}

\begin{tabularx}{\textwidth}{@{}p{0.44\textwidth}Y@{}}
\toprule
Конечное дискретное & без формы хвоста \\
Сглаженная разметка & близкие по риску шаги \\
Смесь геометрических & ранняя остановка и режимы сложности \\
Смесь экспоненциальных & компактный спад хвоста \\
Степенной спад & более медленный хвост \\
Отрицательное биномиальное & тяжелый хвост \\
\bottomrule
\end{tabularx}
\end{column}
\begin{column}{0.46\textwidth}
\textbf{Правила остановки}
\vspace{1mm}

\begin{tabularx}{\textwidth}{@{}p{0.44\textwidth}Y@{}}
\toprule
Накопленная вероятность & $\sum_{j\le k}\widehat p_{t,k}(j)\ge c$ \\
Малое будущее улучшение & $\sum_{j>k}\widehat p_{t,k}(j)\le \varepsilon$ \\
Условная интенсивность & $\widehat p(k)/\sum_{j\ge k}\widehat p(j)\ge c_h$ \\
Квантиль & $q_\alpha\le k$ \\
Бюджет & $\mathbb E\tau\le C$ \\
\bottomrule
\end{tabularx}
\end{column}
\end{columns}
\end{frame}

% 7
\begin{frame}{6. Как получены численные значения}
\begin{enumerate}\setlength\itemsep{2.3mm}
\item Для каждого объекта выполнялись все внутренние шаги до $K=8$.
\item После этого вычислялись потери по шагам и оракульное время $\tau^\star$.
\item Для каждого семейства распределений обучалась голова времени остановки.
\item На контрольной выборке перебирались правила остановки и пороги.
\item В таблицу попадала лучшая строка для данного семейства по контрольному протоколу.
\end{enumerate}

\vspace{2mm}
\smallnote{Метрики: $\bar\tau$ — среднее число шагов; $\textit{tokbest}$ — точность; $NLL_\tau$ —
качество прогноза времени.}
\end{frame}

% 8
\begin{frame}{7. WikiText-103: таблица результатов}
\small
\renewcommand{\arraystretch}{1.18}
\begin{tabularx}{\textwidth}{@{}Y C{14mm} C{15mm} C{16mm} C{15mm}@{}}
\toprule
\textbf{Семейство} & \textbf{Порог} & $\boldsymbol{\bar\tau}$ & \textbf{tokbest} &
$\boldsymbol{NLL_\tau}$ \\
\midrule
\PresWikiTabBody
\bottomrule
\end{tabularx}

\vspace{2mm}
\smallnote{Пять строк --- как в диссертации (нет отдельного прогона \texttt{negative\_binomial} на
WikiText).}

\vspace{2mm}
\begin{block}{Вывод}
\PresWikiTabConclusion
\end{block}
\end{frame}

% 9
\begin{frame}{8. WikiText-103: качество--стоимость}
\centering
{\color{msuBlue}\bfseries Чем левее и выше точка, тем лучше}\par\vspace{1mm}
\begin{tikzpicture}
\begin{axis}[
    scale only axis,
    width=0.82\paperwidth,
    height=0.58\paperheight,
    xmin=\PresWikiQcXmin,
    xmax=\PresWikiQcXmax,
    ymin=\PresWikiQcYmin,
    ymax=\PresWikiQcYmax,
    xlabel={Среднее число внутренних шагов $\bar\tau$},
    ylabel={Top-1 точность},
    grid=both,
    grid style={lineGray!60},
    legend style={font=\footnotesize},
    legend pos=outer north east,
    legend cell align=left,
    tick label style={font=\small},
    label style={font=\small},
    scaled ticks=false,
    yticklabel style={/pgf/number format/fixed,/pgf/number format/zerofill,/pgf/number
    format/precision=3},
    every axis plot/.append style={only marks, mark=*, mark size=3.1pt, draw=pointBlue,
    fill=pointBlue},
]
\PresWikiScatter
\end{axis}
\end{tikzpicture}
\end{frame}

% 10
\begin{frame}{9. ARC-AGI: таблица результатов}
\small
\renewcommand{\arraystretch}{1.18}
\begin{tabularx}{\textwidth}{@{}Y C{14mm} C{15mm} C{16mm} C{15mm}@{}}
\toprule
\textbf{Семейство} & \textbf{Порог} & $\boldsymbol{\bar\tau}$ & \textbf{tokbest} &
$\boldsymbol{NLL_\tau}$ \\
\midrule
\PresArcTabBody
\bottomrule
\end{tabularx}

\vspace{3mm}
\begin{block}{Вывод}
\PresArcTabConclusion
\end{block}
\end{frame}

% 11
\begin{frame}{10. ARC-AGI: качество--стоимость}
\centering
{\color{msuBlue}\bfseries Чем левее и выше точка, тем лучше}\par\vspace{1mm}
\begin{tikzpicture}
\begin{axis}[
    scale only axis,
    width=0.82\paperwidth,
    height=0.58\paperheight,
    xmin=\PresArcQcXmin,
    xmax=\PresArcQcXmax,
    ymin=\PresArcQcYmin,
    ymax=\PresArcQcYmax,
    ytick={0.730,0.735,0.740,0.745},
    xlabel={Среднее число внутренних шагов $\bar\tau$},
    ylabel={Токеновая точность},
    grid=both,
    grid style={lineGray!60},
    legend style={font=\footnotesize},
    legend pos=outer north east,
    legend cell align=left,
    tick label style={font=\small},
    label style={font=\small},
    scaled ticks=false,
    yticklabel style={/pgf/number format/fixed,/pgf/number format/zerofill,/pgf/number
    format/precision=3},
    every axis plot/.append style={only marks, mark=*, mark size=3.1pt, draw=pointBlue,
    fill=pointBlue},
]
\PresArcScatter
\end{axis}
\end{tikzpicture}
\end{frame}

% 12
\begin{frame}{11. Дообучение головы оракула}
\small
\renewcommand{\arraystretch}{1.15}
\begin{tabularx}{\textwidth}{@{}Y C{17mm} C{17mm} C{17mm} C{17mm}@{}}
\toprule
\multirow{2}{*}{\textbf{Семейство}} & \multicolumn{2}{c}{\textbf{WikiText-103: $\bar\tau$}} &
\multicolumn{2}{c}{\textbf{ARC-AGI: $\bar\tau$}} \\
\cmidrule(lr){2-3}\cmidrule(lr){4-5}
& до & после & до & после \\
\midrule
\PresFTtauBody
\bottomrule
\end{tabularx}

\vspace{2mm}
\smallnote{Примечание: для части строк в столбце «до» по WikiText при отсутствии полной базовой
сводки на контроле могут использоваться значения из последнего полного прохода обучения.}
\vspace{2mm}
\begin{block}{Вывод}
Дообучение головы оракула снижает среднюю глубину; на ARC часто сохраняется высокое качество по
тем же таблицам сравнения, что и в основном тексте.
\end{block}
\end{frame}

% 13
\begin{frame}{12. Итог}
\begin{block}{Результат работы}
Сначала оценивается распределение оракульного времени, затем по нему выбирается реализуемый шаг
остановки.
\end{block}

\vspace{2mm}
\begin{columns}[T,onlytextwidth]
\begin{column}{0.48\textwidth}
\textbf{Получено}
\begin{itemize}\setlength\itemsep{1.5mm}
\item формальная разметка оракула через риск;
\item шесть моделей распределения времени;
\item пять правил остановки;
\item сводные эксперименты на двух задачах.
\end{itemize}
\end{column}
\begin{column}{0.48\textwidth}
\textbf{Главный вывод}
\begin{itemize}\setlength\itemsep{1.5mm}
\item смесь геометрических часто дает лучшую точность;
\item смесь экспоненциальных дает сильный компромисс качество--стоимость;
\item конечное дискретное распределение — простая и устойчивая базовая модель.
\end{itemize}
\end{column}
\end{columns}
\end{frame}

\end{document}
